\scanpage{42}
\mspage{41}
\mspage{-3}
6- Coefficients locaux et ``isomorphisme de Thom''.

7- Suites exactes de catégories de Banach.

8-9- Description des groupes $\unsure{k|K}^{p,q}(X)$ à l'aide d'opérateurs de Fredholm. Structures multiplicatives. Espaces classifiants.

10- K-théorie équivariante.

11- Applications diverses de la K-théorie.

N.B.- En principe, aucune connaissance spéciale de K-théorie n'est requise pour suivre cette série de conférences. Cependant, on pourra utilement consulter les exposés de K-théorie du Séminaire Cartan 1963/1964, et les notes aux comptes-rendus, t. 263, série A, p. 275, 347 et 357 et t. 267, même série, p. 305, 328 et 345. (voir aussi l'exposé de Cartan au séminaire Bourbaki 1968/1969).

\uline{Professeur Joseph J. Kohn}: \textit{Complexes différentiels}.

\uline{Programme}

The lectures will concentrate on differential complexes which are not elliptic. The main example of interest is the case of the Cauchy-Riemann equations (the \quad- complex) restricted to a real submanifold. The course will be planned to complement the courses given by Kuranishi and Spencer. If time permits various boundary value problems will also be discussed. Every effort will be made to make the course accessible to all who have had a basic graduate course in analysis.

\margin{$\ldots/$}
